\chapter{The Independence Boundary}
\label{app:independence-boundary}

The main thesis concerns the regular weak null (I(P)=I(Q)>0). This appendix
explains why constructing an independent reference distribution does not turn
the expanded Welch method into a new first-order test of independence.

\section{Constructing the independent reference}

For a joint distribution $P=(p_{ij})$, define
\begin{equation}
  q_{ij}=p_{i+}p_{+j}.
  \label{eq:appendix-independent-q}
\end{equation}
Then $Q=(q_{ij})$ is a probability distribution because
\begin{equation}
  \sum_{i,j}q_{ij}
  =\left(\sum_i p_{i+}\right)
   \left(\sum_j p_{+j}\right)=1.
\end{equation}
It has the same margins as (P), since
\begin{equation}
  \sum_jq_{ij}=p_{i+},
  \qquad
  \sum_iq_{ij}=p_{+j}.
\end{equation}
The factorisation in Equation~\eqref{eq:appendix-independent-q} makes (X)
and (Y) independent under (Q), so (I(Q)=0).

The divergence from (P) to this reference is
\begin{align}
  D_{\mathrm{KL}}(P\Vert Q)
  &=\sum_{i,j}p_{ij}\log\frac{p_{ij}}{q_{ij}}\\
  &=\sum_{i,j}p_{ij}\log\frac{p_{ij}}{p_{i+}p_{+j}}\\
  &=I(P).
\end{align}
Consequently, testing (I(P)=I(Q)) for this constructed reference is exactly
testing (I(P)=0), or independence.

\section{Why the first-order term vanishes}

At independence, $p_{ij}=p_{i+}p_{+j}$, and every population pointwise-MI
value is zero:
\begin{equation}
  \pmi_P(i,j)=\log 1=0.
\end{equation}
Equation~\eqref{eq:mi-derivative} therefore gives
\begin{equation}
  \left.\frac{\mathrm d I(P_\varepsilon)}
  {\mathrm d\varepsilon}\right|_0
  =\pmi_P(x,y)-I(P)=0.
\end{equation}
The first-order variance also vanishes:
\begin{equation}
  V(P)=\Var_P\{\pmi_P(X,Y)\}=0.
\end{equation}
The denominator used by normal Wald, simple Welch, and expanded Welch is thus
degenerate under the population independence null. Changing its degrees of
freedom cannot repair a statistic whose leading term is zero.

\section{Second-order behaviour}

Let an empirical perturbation around an independent distribution have
first-order cell residuals $h_{ij}$, with row and column sums $h_{i+}$ and
$h_{+j}$. The association component is
\begin{equation}
  a_{ij}=h_{ij}-p_{+j}h_{i+}-p_{i+}h_{+j}.
\end{equation}
The linear contribution to MI is zero. Expanding the logarithm to second
order gives
\begin{equation}
  I(P_\varepsilon)
  =\frac{\varepsilon^2}{2}
  \sum_{i,j}\frac{a_{ij}^2}{p_{i+}p_{+j}}
  +o(\varepsilon^2).
  \label{eq:appendix-mi-second-order}
\end{equation}
For multinomial sampling, empirical probability errors are of order
$n^{-1/2}$. Equation~\eqref{eq:appendix-mi-second-order} is therefore of
order $n^{-1}$, and the nondegenerate statistic is $2n\Ihat(P)$, not
$\sqrt n\{\Ihat(P)-I(P)\}$. Its regular fixed-alphabet limit is
\begin{equation}
  2n\Ihat(P)
  \xrightarrow{d}
  \chi^2_{(r-1)(c-1)}.
\end{equation}
This is the usual likelihood-ratio test of independence. The chi-squared
reference arises from the second-order quadratic geometry at the boundary,
whereas the expanded Welch reference arises from first-order variation away
from that boundary.

\section{What alternative constructions produce}

Using the observed margins to set
$\widehat q_{ij}=\widehat p_{i+}\widehat p_{+j}$ gives
\begin{equation}
  2nD_{\mathrm{KL}}(\widehat P\Vert\widehat Q)
  =2n\Ihat(P),
\end{equation}
which is the likelihood-ratio statistic above.

Repeatedly shuffling one variable preserves its observations while breaking
the pairing with the other variable. The resulting empirical distribution is
a permutation reference. Repeatedly drawing tables from
$\widehat P_X\widehat P_Y$ is a parametric bootstrap reference. A single
simulated $Q$ table adds Monte Carlo variation but does not by itself define
a calibrated null distribution.

The independent reference $Q=P_XP_Y$ is therefore mathematically valid, but
it does not create a regular two-sample MI comparison. Analytic improvement of
the sparse independence test would require a second-order finite-sample
method, rather than an immediate extension of the expanded Welch calculation.
